% ======================================================================
%  sec08.tex —— 第 8 节 点函子
%  底本：原书 pp.44–45
%  原书标题：8. Functor of points
%            8.1 R-points / 8.2 Schemes as a functor
% ======================================================================

% 原书的混合环境在导言区未定义，此处就地补上（aliascnt 使其与定理共享计数器）。
\makeatletter
\@ifundefined{c@propdef}{\newaliascnt{propdef}{theorem}\newtheorem{propdef}[propdef]{命题--定义}\aliascntresetthe{propdef}}{}
\@ifundefined{c@propcon}{\newaliascnt{propcon}{theorem}\newtheorem{propcon}[propcon]{命题--构造}\aliascntresetthe{propcon}}{}
\makeatother

\section{点函子}\label{sec:8}

本节的主要目标是构造完全忠实的函子 $\Sch\to\mathsf{Fun}(\CRing,\Set)$，并描述它的本质像。

\subsection{\texorpdfstring{$R$}{R}-点}\label{sec:8.1}

\begin{definition}\label{def:8.1.1}
设 $X$ 为概形，$R$ 为交换环。$X$ 的一个\emph{$R$-点}是指态射 $\Spec(R)\to X$。令
\[
X(R):=\Hom_{\Sch}\bigl(\Spec(R),X\bigr)
\]
为 $X$ 的 $R$-点集。

$X$ 的\emph{域值点}是指 $X$ 的某个域 $k$ 上的 $k$-点。$X$ 的\emph{几何点}是指 $X$ 的某个可分闭域 $k$ 上的 $k$-点。
\end{definition}

\begin{propcon}\label{prop:8.1.2}
设 $X$ 为概形，$k$ 为域。则下述两个集合之间存在双射：
\begin{itemize}[leftmargin=1.6em,itemsep=0.2ex]
  \item $X$ 的 $k$-点集 $X(k)$；
  \item 二元组 $(x,i)$ 构成的集合，其中 $x\in X$ 是拓扑点\footnote{指底层拓扑空间中的点。}，$i:\kappa_x\to k$ 是环同态。
\end{itemize}
\end{propcon}

\begin{proof}
按定义，$X$ 的 $k$-点由二元组 $(\varphi,\beta)$ 给出。由于 $\Spec(k)$ 只有一个拓扑点，连续映射 $\varphi$ 由某个拓扑点 $x\in X$ 给出。此时态射 $\beta$ 由局部同态 $\cO_{X,x}\to k$ 给出；这里把单点空间 $\{*\}$ 上交换环层 $\mathcal F$ 与 $\mathcal F(\{*\})$ 等同。而局部同态 $\cO_{X,x}\to k$ 唯一地分解为
\[
\cO_{X,x}\to\kappa_x\to k .
\]
这就给出所要的双射。
\end{proof}

\begin{remark}\label{rem:8.1.3}
对域值点 $s:\Spec(k)\to X$，有时滥用记号写作 $s\to X$，其中 $s$ 理解为 $\Spec(k)$。类似地，对拓扑点 $x\in X$，有时滥用记号把 $\kappa_x$-点 $\Spec(\kappa_x)\to X$ 写作 $x\to X$。

但要注意，概形并不被它的域值点决定。
\end{remark}

\begin{exercise}\label{exc:8.1.4}
设 $k$ 为域，考虑 $X:=\Spec(k)$ 与 $X':=\Spec(k[\epsilon]/(\epsilon^2))$。设 $f:X\to X'$ 为对应显然同态 $k[\epsilon]/(\epsilon^2)\to k$ 的态射。证明 $f$ 在 $X$ 与 $X'$ 的域值点集之间诱导双射。
\end{exercise}

\begin{construction}\label{con:8.1.5}
设 $X$ 为概形，$x\in X$ 为拓扑点。取含 $x$ 的仿射开子集 $U\subseteq X$，注意 $U\simeq\Spec(\cO_X(U))$。典范同态 $\cO_X(U)\to\cO_{X,x}$ 诱导态射 $\Spec(\cO_{X,x})\to\Spec(\cO_X(U))$。于是合成
\[
\Spec(\cO_{X,x})\to\Spec(\cO_X(U))\simeq U\to X
\]
给出 $X$ 的一个 $\cO_{X,x}$-点。
\end{construction}

\begin{exercise}\label{exc:8.1.6}
证明：
\begin{enumerate}[label=(\arabic*),leftmargin=2em,itemsep=0.25ex]
  \item 上述 $X$ 的 $\cO_{X,x}$-点不依赖 $U$ 的选取；
  \item 态射 $\Spec(\cO_{X,x})\to X$ 把局部环化空间 $\Spec(\cO_{X,x})$ 实现为 $X$ 限制到某个子空间所得的局部环化空间。
\end{enumerate}
\end{exercise}

\subsection{概形作为函子}\label{sec:8.2}

\begin{definition}\label{def:8.2.1}
设 $X$ 为概形。函子
\[
h_X:\CRing\to\Set,\qquad R\mapsto X(R)
\]
称为 $X$ 的\emph{点函子}。

注意构造 $X\mapsto h_X$ 具有函子性。于是得到典范函子
\[
h:\Sch\to\mathsf{Fun}(\CRing,\Set).
\]
\end{definition}

\begin{remark}\label{rem:8.2.2}
按构造，合成
\[
\CRing^{\mathrm{op}}\simeq\mathsf{Aff}\to\Sch\xrightarrow{\ h\ }\mathsf{Fun}(\CRing,\Set)
\]
可等同于 Yoneda 嵌入 $R\mapsto\Hom_{\CRing}(R,-)$。
\end{remark}

\begin{theorem}\label{thm:8.2.3}
函子 $h:\Sch\to\mathsf{Fun}(\CRing,\Set)$ 完全忠实。
\end{theorem}

\begin{remark}\label{rem:8.2.4}
关于该函子本质像的描述\footnote{在标准术语下，函子 $F:\CRing\to\Set$ 属于本质像，当且仅当 $F$ 满足 Zariski 下降，且可被可表函子的开覆盖覆盖。}，参见 \cite[Chapter VI]{EH00}。
\end{remark}

\begin{remark}\label{rem:8.2.5}
上述定理提示了另一条发展概形理论的途径，它不使用局部环化空间。具体地说，可以把概形定义为满足某些性质的函子 $\CRing\to\Set$。事实上，在 1970 年代，Grothendieck 本人就极力主张放弃他早期给出的概形定义，转而采用函子观点。在笔者看来，这一途径至少有如下优点。
\begin{itemize}[leftmargin=1.6em,itemsep=0.2ex]
  \item 它使大量关于概形的构造变得形式化，因而更容易。
  \item 它提供了处理模问题与形变理论的更直接的途径。
  \item 它允许定义更奇异的几何对象，例如代数空间、叠、indscheme 等。
\end{itemize}
然而，无论笔者多么喜爱这种函子途径，都不认为初学者可以忽视「概形是带结构的拓扑空间」这一经典观点\footnote{试想第一次学习射影空间 $\PP^n_\Z$ 就用下述定义：它是把交换代数 $R$ 送到「自由 $R$-模 $R^{\oplus(n+1)}$ 到秩 $1$ 投射 $R$-模 $P$ 的满射的同构类」的函子。}。
\end{remark}

\begin{definition}\label{def:8.2.6}
称函子 $F:\CRing\to\Set$ 被概形 $X$ \emph{表示}，若它带自然同构 $F\simeq h_X$。
\end{definition}

\begin{example}\label{ex:8.2.7}
把 $R$ 送到其底层集合的函子 $\CRing\to\Set$ 由仿射概形 $\A^1_\Z:=\Spec(\Z[t])$ 表示。
\end{example}

\begin{exercise}\label{exc:8.2.8}
证明：把 $R$ 送到 $R$ 上 $n\times n$ 可逆矩阵集 $\mathrm{GL}_n(R)$ 的函子 $\CRing\to\Set$ 由某个仿射概形表示。
\end{exercise}

\begin{exercise}\label{exc:8.2.9}
证明：常函子 $\CRing\to\Set$，$R\mapsto I$ 不被概形表示，除非 $I\simeq\{*\}$。不交并 $\coprod_{i\in I}\Spec(\Z)$ 表示的是什么函子？
\end{exercise}

\begin{yizhu}
$R$-点集 $X(R)=\Hom_{\Sch}(\Spec(R),X)$ 对 $R$ 是反变的：环同态 $R\to R'$ 诱导 $\Spec(R')\to\Spec(R)$，从而给出 $X(R)\to X(R')$。这正是把 $X$ 写成函子 $h_X:\CRing\to\Set$ 的全部内容，也是「点函子」一词的由来——对每个测试环 $R$，$X(R)$ 就是 $X$ 的「$R$ 值点」。
[来源：Stacks Project, Tag 01JF（Section 26.15, A representability criterion），\url{https://stacks.math.columbia.edu/tag/01JF}（2026-11-17 访问）。]
\end{yizhu}

\begin{yizhu}
命题--构造 \ref{prop:8.1.2} 是 Stacks 中一条更一般引理的特例：对局部环 $R$，态射 $\Spec(R)\to X$ 与二元组 $(x,\ \cO_{X,x}\to R)$ 一一对应，其中第二个分量是局部同态。把 $R$ 取成域 $k$，局部同态只能经由剩余域 $\kappa_x$ 分解，于是得到本题的双射。
[来源：Stacks Project, Tag 01J6（Lemma 26.13.1, Points of schemes），\url{https://stacks.math.columbia.edu/tag/01J6}（2026-11-17 访问）。]
\end{yizhu}

\begin{yizhu}
习题 \ref{exc:8.1.4} 的现象值得记住：$\Spec(k[\epsilon]/(\epsilon^2))$ 与 $\Spec(k)$ 有相同的域值点，却带着一层非平凡的幂零信息。这说明域值点看不见幂零元，因而看不见非既约的结构；要看见它们，必须允许把 $\Spec(k[\epsilon]/(\epsilon^2))$ 这样的「加厚点」当作测试对象。
[依据：原书 p.44 习题 8.1.4 与注记 8.1.3]
\end{yizhu}

\begin{yizhu}
定理 \ref{thm:8.2.3} 的完全忠实性可以看作定理 \ref{thm:7.1.2} 的函子版：仿射情形给出 $\Hom(\Spec R,\Spec R')\cong\Hom(R',R)$，而一般概形由仿射开集粘接而成，粘接数据在函子层面恰好转化为下降数据。原书没有给出证明，但这条定理是后面用 Yoneda 引理处理概形问题的合法依据。
[依据：原书 p.45 定理 8.2.3 与 p.42 定理 7.1.2]
\end{yizhu}

\begin{yizhu}
例 \ref{ex:8.2.7} 之所以成立，是因为 $\Hom_{\CRing}(\Z[t],R)$ 与 $R$ 的元素一一对应：给定环同态，$t$ 的像可任取；反之任取 $r\in R$，由 $\Z[t]\to R$，$t\mapsto r$ 得到唯一同态。于是 $\A^1_\Z$ 表示「忘记环结构、只留底层集合」的函子。
[译者按（上下文推断，无外部依据）：由自由多项式环的泛性质直接验证。]
\end{yizhu}

\begin{yizhu}
习题 \ref{exc:8.2.9} 的答案可由命题 \ref{prop:6.1.3} 读出：不交并是余积，故
\[
\Hom_{\Sch}\Bigl(\Spec(R),\coprod_{i\in I}\Spec(\Z)\Bigr)\cong\prod_{i\in I}\Hom_{\Sch}\bigl(\Spec(R),\Spec(\Z)\bigr),
\]
而每个因子只有一个元素，因为 $\Z\to R$ 的环同态唯一。于是当 $I$ 非空时，$\coprod_{i\in I}\Spec(\Z)$ 表示取值为单点的常函子。
[译者按（上下文推断，无外部依据）：用不交并的余积性质与 $\Z$ 的初始性。]
\end{yizhu}
